\documentclass[12pt]{article}
\usepackage[margin=0.58in,top=0.62in,bottom=0.48in]{geometry}
\usepackage{lmodern}
\usepackage{microtype}
\usepackage{fancyhdr}
\usepackage{titlesec}
\usepackage{enumitem}
\usepackage{lastpage}
\usepackage[hidelinks]{hyperref}
\setlength{\parindent}{0pt}
\setlength{\parskip}{0.18em}
\setlength{\headheight}{26pt}
\raggedbottom
\titleformat{\section}{\normalsize\bfseries\scshape}{}{0pt}{}
\titleformat{\subsection}{\normalsize\bfseries}{}{0pt}{}
\titlespacing*{\section}{0pt}{0.35em}{0.12em}
\titlespacing*{\subsection}{0pt}{0.25em}{0.08em}
\pagestyle{fancy}
\fancyhf{}
\fancyhead[C]{\footnotesize Lit Example Reader \quad Assignment ID: U1L19LE \quad Page \thepage\ of \pageref{LastPage}}
\renewcommand{\headrulewidth}{0.5pt}
\newenvironment{playpassage}{\begin{quote}\footnotesize\setlength{\parskip}{0.08em}\setlength{\parindent}{0pt}}{\end{quote}}
\newcommand{\speaker}[1]{\textbf{#1.}}
\begin{document}
\begin{center}
{\LARGE\bfseries Week 19 Lit Example Reader}\par\vspace{0.02em}
{\large\scshape Tragic Judgment and Disciplined Thought}\par\vspace{0.06em}
\rule{0.78\textwidth}{0.5pt}
\end{center}

\textbf{Purpose:} Use Shakespeare as a judgment mirror in the final portfolio defense: analyze a failure precisely, then connect the warning to a disciplined habit.

\textbf{What to watch for:} Do not claim moral superiority over a character. Identify a warning, the inference or judgment made under pressure, and the course habit that tests it.

\section*{Before the Passages: Literary Evidence in a Defense}
\begin{itemize}[leftmargin=1.3em,itemsep=0.05em,topsep=0.08em,parsep=0.03em]
\item \textbf{Claim:} a precise judgment about the character's reasoning.
\item \textbf{Evidence:} exact words and dramatic circumstances.
\item \textbf{Warrant:} why the evidence demonstrates a disciplined or undisciplined habit.
\item \textbf{Transfer:} a concrete setting where the warning changes how a thinker acts.
\end{itemize}

\section*{Passage 1: Macbeth, Act 5, Scene 5 --- Confidence Meets Discovery}
\textbf{Context:} A messenger reports that Birnam Wood appears to move. Macbeth first attacks the report, then recognizes that the apparition's reassuring words may have trapped him.
\begin{playpassage}
\speaker{Messenger} As I did stand my watch upon the hill,\\
I look'd toward Birnam, and anon, methought,\\
The wood began to move.\\[0.2em]
\speaker{Macbeth} Liar and slave!\\
...\\
I pull in resolution, and begin\\
To doubt the equivocation of the fiend\\
That lies like truth: ``Fear not, till Birnam wood\\
Do come to Dunsinane''; and now a wood\\
Comes toward Dunsinane.
\end{playpassage}

\subsection*{Reader Notes}
This is the assisted example. The new observation conflicts with Macbeth's settled confidence. His first response attacks the witness; only then does he revisit the source and ambiguous promise. A final defense could use this moment to support a limited claim: disciplined judgment remains open to disconfirming evidence and retests source and terms when facts resist expectation. The passage does not prove that one habit prevents every disaster; it gives a vivid warrant for intellectual humility.

\textbf{Reading task:} Underline the observation, box the first reaction, and star the line where Macbeth begins to revise his judgment. In the margin, state the habit this sequence warns a thinker to practice.

\vfill
\section*{Annotation Notes}
\noindent\rule{\textwidth}{0.4pt}\\[1.0em]
\rule{\textwidth}{0.4pt}\\[1.0em]
\rule{\textwidth}{0.4pt}

\newpage
\section*{Passage 2: Julius Caesar, Act 2, Scene 1 --- A Decision Before the Evidence}
\textbf{Context:} Brutus turns a possible future change in Caesar into a present conclusion that death is necessary.
\begin{playpassage}
\speaker{Brutus} It must be by his death: and for my part,\\
I know no personal cause to spurn at him,\\
But for the general. He would be crown'd:\\
How that might change his nature, there's the question.\\
...\\
And therefore think him as a serpent's egg\\
Which, hatch'd, would, as his kind, grow mischievous,\\
And kill him in the shell.
\end{playpassage}

\subsection*{Reader Notes}
This passage requires independent transfer. Attend to the movement from possibility to necessity, the analogy that carries the movement, and the absent alternatives. Decide which one or two Whately habits best expose the pressure point. Then ask what a disciplined thinker would do before accepting the conclusion. Do not turn the passage into a list of fallacy names or a plot summary.

\textbf{Reading task:} Underline the possibility, box the conclusion, star the analogy, and write one fair counterpressure the conclusion must answer.

\section*{How to Use These Passages on the Worksheet}
Build a miniature defense: make a claim about disciplined judgment, use one passage as evidence, state the warrant, acknowledge a limit, and name a concrete transfer. Passage 1 offers a model path; Passage 2 asks you to choose and defend the logical tool.

\vfill
\section*{Defense Notes}
\noindent\rule{\textwidth}{0.4pt}\\[1.0em]
\rule{\textwidth}{0.4pt}\\[1.0em]
\rule{\textwidth}{0.4pt}\\[1.0em]
\rule{\textwidth}{0.4pt}\\[1.0em]
\rule{\textwidth}{0.4pt}
\end{document}
